\subsection{Описание алгоритма Дейкстры}

Алгоритм состоит из двух основных компонентов: \textbf{генерация графа} и \textbf{алгоритм Дейкстры}.

Функция \texttt{generateGraph(int n, int m, int q, int r)} генерирует случайный неориентированный граф с \texttt{n} вершинами и до \texttt{m} ребер. Вес ребер присваиваются случайным образом в диапазоне от \texttt{q} до \texttt{r}.

Максимальное количество ребер ограничено формулой \( \frac{n(n-1)}{2} \), которая представляет количество ребер в полном графе. Двумерный вектор \texttt{was} используется для отслеживания существующих ребер, чтобы избежать дублирования. Ребра добавляются в список смежности \texttt{adj\_list} до достижения желаемого количества ребер.

\begin{lstlisting}
m = min(n * (n - 1) / 2, m);
vector<vector<pair<int, int>>> adj_list(n);
std::random_device rd;
std::mt19937 gen(rd());
std::uniform_int_distribution<> dist(q, r);
vector<vector<bool>> was(n, vector<bool>(n, false));
int edge_count = 0;

while (edge_count < m) {
    int i = rand() % n;
    int j = rand() % n;
    if (i != j && !was[i][j]) {
        int weight = dist(gen);
        adj_list[i].push_back({ j, weight });
        adj_list[j].push_back({ i, weight });
        was[i][j] = was[j][i] = true;
        ++edge_count;
    }
}
\end{lstlisting}

Функция реализует алгоритм Дейкстры для нахождения кратчайшего пути от стартовой вершины ко всем другим вершинам графа.

Алгоритм инициализирует расстояния до бесконечности (\texttt{INT\_MAX}) и устанавливает расстояние до стартовой вершины равным нулю. Приоритетная очередь, представляемая D-ариHeap (\texttt{DHeap}), используется для эффективного извлечения вершины с минимальным расстоянием. Алгоритм итеративно обрабатывает вершины, обновляя расстояния до их соседей, если найден более короткий путь.

\begin{lstlisting}
vector<bool> processed(n, false);
vector<int> index(n);
vector<int> name(n);

for (int i = 0; i < n; ++i) {
    up[i] = -1;
    dist[i] = INT_MAX;
    index[i] = i;
    name[i] = i;
}

dist[start] = 0;
up[start] = start;
DHeap dHeap(d);
dHeap.Insert(start, dist[start]);

while (!dHeap.IsEmpty()) {
    pair<int, int> minPair = dHeap.ExtractMin();
    int i = minPair.first;
    int currentDist = minPair.second;

    if (processed[i]) continue; 
    processed[i] = true;

    vector<pair<int, int>> edge = graph[i];
    for (int k = 0; k < edge.size(); k++) {
        int j = edge[k].first;
        int newDist = edge[k].second + currentDist;

        if (newDist < dist[j]) {
            dist[j] = newDist;
            up[j] = i;
            dHeap.Insert(j, dist[j]);
        }
    }
}
\end{lstlisting}

Время выполнения алгоритма Дейкстры измеряется с помощью высокоточных часов из библиотеки \texttt{chrono}. Время, затраченное на выполнение алгоритма, выводится в конце.

